% !TEX root = ../Thesis.tex

\chapter{Fazit und Ausblick}
\label{ch:abschluss}%

\section{Zusammenfassung der Ergebnisse}

Die Arbeit untersuchte die Entwicklung und Evaluation\index{Evaluation!Agenten} agentischer Architekturen\index{Referenzarchitektur} für Software-Engineering-Workflows\index{SE-Workflows}.\footnote{Vergleiche dazu~\cite{xi2023rise} und~\cite{wang2023survey} für aktuelle Forschungstrends.}
Ausgehend von einer systematischen Analyse des Stands der Forschung (Kapitel~\ref{ch:hintergrund}) wurde eine Referenzarchitektur entwickelt\index{Architekturdesign} (Kapitel~\ref{ch:konzept}), prototypisch implementiert und empirisch evaluiert\index{Empirische Evaluation} (Kapitel~\ref{ch:realisierung}).

Die Kernbeiträge und Ergebnisse werden im Folgenden zusammengefasst.

\subsection{Beantwortung der Forschungsfragen}

\textbf{Forschungsfrage 1:} Wie lassen sich robuste Agenten-Policies\index{Agenten-Policies!Robustheit} für SE-Workflows systematisch modellieren\index{Policy-Modellierung}?

Durch Kombination des ReAct-Patterns\index{ReAct!Anwendung} (Reasoning + Acting) mit expliziten Reflexionsmechanismen\index{Reflexionsmechanismen} konnten strukturierte Strategien\index{Strukturierte Policies} entwickelt werden, die Planung\index{Planung!in Policies}, Werkzeugorchestrierung und Selbstkritik integrieren. Die Evaluation zeigte, dass die Strategien Erfolgsraten von \pct{73} bei Refactoring-Aufgaben erreichen \textendash{} signifikant über Baseline-Systemen (\pct{42}--\pct{58}).

Zentrale Design-Entscheidungen umfassten die Implementierung expliziter Denk-Handlungs-Beobachtungs-Schleifen\index{Thought-Action-Observation}, die Transparenz\index{Agent-Transparenz} durch nachvollziehbare Zwischenschritte schaffen. Strukturierte Werkzeugschnittstellen\index{Tool-Interfaces!Strukturierung} mit JSON-Schema-Validierung\index{Schema-Validation} gewährleisten typsichere Kommunikation zwischen Komponenten. Das episodische Gedächtnis\index{Episodisches Gedächtnis!Design} ermöglicht Kontextpersistierung\index{Kontext-Persistierung} über Iterationen hinweg und unterstützt damit langfristige, zustandsabhängige Workflows.

\textbf{Forschungsfrage 2:} Welche Architekturprinzipien ermöglichen sichere\index{Sicherheit!Architektur} und nachvollziehbare\index{Nachvollziehbarkeit} Tool-Integration\index{Tool-Integration!Sicherheit}?

Die entwickelte Architektur implementiert Verteidigung in der Tiefe (Defense-in-Depth)\index{Defense-in-Depth} durch mehrschichtige Sicherheitsmechanismen. Alle LLM-generierten Werkzeugaufrufe unterliegen einer Eingabevalidierung\index{Input-Validation!Tool-Calls}, die schädliche Aufrufe verhindert. Die Sandbox-Ausführung\index{Sandbox!Docker} in isolierten Docker-Containern stellt sicher, dass potenziell gefährlicher Code nicht das Host-System kompromittieren kann. Ein Berechtigungssystem\index{Permissions-System} setzt das Prinzip der geringsten Rechte (Least-Privilege)\index{Least-Privilege} durch und beschränkt Zugriffe auf das erforderliche Minimum. Kritische Operationen werden durch Audit-Protokollierung\index{Audit-Logging!Tool-Calls} lückenlos protokolliert. Bei deploymentkritischen Aktionen greift ein Mechanismus zur menschlichen Freigabe, der eine explizite Bestätigung verlangt.

Sicherheits-Audits\index{Sicherheits-Audits} zeigten \pct{100} Erfolgsrate bei der Abwehr von Prompt-Injection-Angriffen\index{Prompt-Injection} und unautorisiertem Dateizugriff\index{File-Access-Kontrolle}.

\textbf{Forschungsfrage 3:} Mit welchen Metriken kann die Leistungsfähigkeit agentischer Systeme realistisch bewertet werden?

Ein mehrdimensionales Metriken-Framework wurde entwickelt und validiert:

\begin{description}
  \item[Funktional:] Aufgabenerfolgsquote, Testbestehensquote, Lint-Bewertung, Review-Akzeptanz
  \item[Effizienz:] Token-Verbrauch, API-Kosten, Laufzeit, Anzahl Werkzeugaufrufe
  \item[Robustheit:] Fehlerbehebungsrate, Graceful Degradation bei Failures
  \item[Sicherheit:] Policy-Verstöße, Sandbox-Ausbrüche, Zugriffsverweigerungen
\end{description}

Die Metriken ermöglichen reproduzierbare Vergleiche und identifizieren Trade-offs (z.\,B. Reflexion erhöht Erfolgsrate um \pct{15}, aber Kosten um \pct{12}).

\subsection{Empirische Validierung}

Die prototypische Implementierung wurde anhand von 25 realistischen SE-Szenarien evaluiert:

\begin{itemize}
  \item \textbf{Refactoring:} \pct{73} Erfolgsrate (Extract-Function, Rename, Simplify)
  \item \textbf{Test-Debugging:} \pct{62} Erfolgsrate (Diagnose + Fix)
  \item \textbf{Lint-Resolution:} \pct{86} Erfolgsrate (Style + Type Errors)
  \item \textbf{Durchschn. Kosten:} \$0.08 pro erfolgreichem Task
  \item \textbf{Token-Effizienz:} \pct{40} Reduktion vs. naive Baseline durch Kontextoptimierung
  \item \textbf{Rechenumgebungen:} Evaluation sowohl ohne als auch mit \textbf{GPU}-Beschleunigung
\end{itemize}

Vergleiche mit Baseline-Systemen (naives Prompting, Standard-ReAct) zeigten +31\% bzw. +15\% Erfolgsraten-Verbesserungen.

\section{Beiträge der Arbeit}

Die Arbeit leistet folgende Beiträge zur Spezialisierung \enquote{Software Engineering mit agentischer KI}:

\begin{enumerate}
  \item \textbf{Referenzarchitektur:} Dokumentierte, wiederverwendbare Architektur für agentische SE-Workflows mit klaren Komponenten (Controller, Werkzeugregister, Gedächtnismanager, Sicherheitsschicht) und deren Interaktionen. Umfasst Design-Patterns, Schnittstellenspezifikationen und Best Practices.
  
  \item \textbf{Praxisnahe Implementierung:} Open-Source-Prototyp in Python mit vollständiger Tool-Integration (Linter, Tests, VCS). Inkl. Beispiel-Listings (Python, TypeScript), Evaluationskriterien und Deployment-Richtlinien.
  
  \item \textbf{Sicherheitsframework:} Konkrete Mechanismen für sichere Agentenausführung: Eingabevalidierung, Sandboxing, Berechtigungskonzepte, Audit-Logging und menschliche Freigabe. Getestet gegen Prompt-Injection und unautorisierten Zugriff.
  
  \item \textbf{Evaluationsmethodik:} Mehrdimensionales Metrikenframework (Funktional, Effizienz, Robustheit, Sicherheit) mit reproduzierbaren Benchmarks. Ermöglicht systematische Vergleiche und Trade-off-Analysen.
  
  \item \textbf{Empirische Evidenz:} Die Validierung anhand von 25 realen SE-Szenarien zeigt praktische Durchführbarkeit und quantifiziert Verbesserungen (+31\% Erfolgsrate, -40\% Token-Kosten) gegenüber Baselines.
  
  \item \textbf{Transferierbarkeit:} Architektur ist nicht auf SE beschränkt. Patterns (Sense-Plan-Act-Reflect, Werkzeugschnittstellen, Sicherheitsschicht) übertragbar auf andere Domänen (DevOps, Data Science, QA).
\end{enumerate}

\subsection{Wissenschaftlicher Beitrag}

Im Kontext der Forschungslandschaft (vgl. Kapitel~\ref{ch:hintergrund}) schließt die Arbeit folgende Lücken:

\begin{itemize}
  \item Systematisierung agentischer SE-Architekturen (bisher meist ad-hoc Prototypen)
  \item Fokus auf Produktionsreife (Sicherheit, Kosten, Robustheit) statt nur Aufgabenerfolg
  \item Transferierbare Design-Patterns statt monolithischer Systeme
  \item Vergleichbare Evaluation mit reproduzierbaren Benchmarks
\end{itemize}

\section{Limitierungen}

Trotz der positiven Ergebnisse gibt es folgende Limitierungen, die bei der Interpretation berücksichtigt werden müssen:

\subsection{Methodische Limitierungen}

\begin{itemize}
  \item \textbf{Begrenzte Testmenge:} Die Evaluation auf 25 Tasks (3 Szenarien) liefert belastbare Indikationen, ist aber nicht hinreichend für abschließende Aussagen zur Produktionsreife. Größere und methodisch strengere Evaluationssetups, etwa SWE-bench Pro oder vergleichbare End-to-End-Benchmarks, wären wünschenswert~\cite{openai2026swebenchverified,swebench2026leaderboard}.
  
  \item \textbf{Kontrollierte Umgebung:} Experimente erfolgten auf kuratierten Projekten mit sauber definierten Szenarien. Praxiseinsätze weisen häufig ambigere Anforderungen, Altsysteme und unvollständige Dokumentation auf.
  
  \item \textbf{Skalierungsgrenzen:} Tests beschränkten sich auf Projekte bis 50k LOC. Verhalten bei Millionen LOC (Linux Kernel, Chromium) ist unklar.
  
  \item \textbf{LLM-Abhängigkeit:} Die Ergebnisse basieren auf Frontier-Modellen der GPT- und Claude-Klasse. Neuere Modellgenerationen können Architektur-Trade-offs verschieben; kleinere oder offene Modelle führen je nach Aufgabe zu abweichenden Kosten-, Latenz- und Qualitätsprofilen.
\end{itemize}

\textquote[\cite{jimenez2023swe}]{Während agentische Systeme Potenzial zeigen, müssen ihre Grenzen in kontrollierten Umgebungen sorgfältig evaluiert werden, bevor sie in Produktionssystemen eingesetzt werden.}

\subsection{Technische Limitierungen}

\begin{itemize}
  \item \textbf{Halluzinationen:} LLMs erzeugen gelegentlich nicht existente APIs oder fehlerhafte Begründungsschritte. Reflexion reduziert dieses Risiko, eliminiert es jedoch nicht.
  
  \item \textbf{Kontextfenster:} Trotz deutlicher Fortschritte bei langen Kontextfenstern stoßen selbst Modelle mit sehr großen Eingabebudgets bei komplexen Codebasen an praktische Grenzen~\cite{meta2025llama4}. RAG- und Chunking-Strategien bleiben daher Hilfsmittel, keine vollständige Lösung.
  
  \item \textbf{Werkzeuglatenz:} Testläufe dauern Sekunden bis Minuten. Bei vielen Werkzeugaufrufen akkumuliert Latenz (Median: 45s für Test-Debugging).
  
  \item \textbf{Fehlerfortpflanzung:} Frühe Fehler (z.\,B. falsche Diagnose) propagieren durch iterative Schleifen. Ohne menschliche Aufsicht können Agenten in ineffiziente oder fehlerhafte Suchpfade geraten.
\end{itemize}

\subsection{Gesellschaftliche und ethische Limitierungen}

Zu berücksichtigen ist auch die gesellschaftliche Dimension. Die verfügbare Literatur deutet darauf hin, dass LLMs und agentische Systeme Tätigkeitsprofile im Software Engineering spürbar verändern können, ohne dass sich daraus bereits lineare oder einheitliche Aussagen über Substitutionseffekte ableiten lassen~\cite{eloundou2023gpts}. Für wissenschaftliche und industrielle Anwendungen bedeutet dies, dass Produktivitätsgewinne, Qualifikationsverschiebungen und Verantwortungsfragen gemeinsam betrachtet werden müssen.

Weitere ethische Dimensionen:

\begin{itemize}
  \item \textbf{Bias-Verstetigung:} LLMs reproduzieren Biases aus Trainingsdaten. Code-Generierung kann diskriminierende Patterns fortführen.
  
  \item \textbf{Verantwortlichkeit:} Bei agentengenerierten Fehlern bleibt die Zurechnung zwischen nutzender Organisation, verantwortlicher Fachperson und Modellanbieter eine offene Governance-Frage.
  
  \item \textbf{Übermäßiges Vertrauen:} Entwicklungsteams könnten kritisches Prüfen reduzieren und Agenten-Outputs zu unkritisch übernehmen \textendash{} besonders problematisch in sicherheitskritischen Systemen.
\end{itemize}

\section{Zukünftige Arbeiten}

Auf Basis dieser Arbeit ergeben sich mehrere Richtungen für zukünftige Forschung und Entwicklung:

\subsection{Kurzfristige Erweiterungen}

\begin{itemize}
  \item \textbf{Erweiterte Benchmarks:} Evaluation auf methodisch robusteren Benchmarks wie SWE-bench Pro sowie auf HumanEval-ähnlichen Datensätzen für breitere Validierung
  
  \item \textbf{Mehrsprachenunterstützung:} Aktuell Python-fokussiert. Erweiterung auf Java, TypeScript, Go für breitere Anwendbarkeit
  
  \item \textbf{Optimiertes Kontextmanagement:} Hybrid-Strategien (RAG + Summarization + Code-Graph-Navigation) für sehr große Codebasen mit langen Kontextfenstern
  
  \item \textbf{Fine-Tuning:} Domänenspezifisches Fein-Tuning auf SE-Szenarien könnte Erfolgsraten bei kleineren und günstigeren Modellen verbessern
  
  \item \textbf{Integration menschlichen Feedbacks:} RLHF-ähnliche Ansätze für kontinuierliches Lernen aus Korrekturen im Entwicklungsteam
\end{itemize}

\subsection{Mittel- bis langfristige Forschungsrichtungen}

\begin{itemize}
  \item \textbf{Multi-Agenten-Systeme:} Rollenbasierte Kollaboration wie in MetaGPT. Dies könnte Spezialisierung und Parallelisierung verbessern.
  
  \item \textbf{Kontinuierliches Lernen:} Agenten lernen aus Projekthistorie, Teamkonventionen und Codebasis-Mustern. Episodisches Gedächtnis könnte dabei als Datenquelle dienen.
  
  \item \textbf{Formale Verifikation:} Integration formaler Methoden (Typprüfung, SMT-Solving, symbolische Ausführung) für sicherheitskritischen Code.
  
  \item \textbf{IDE-Integration:} Tiefe Integration in Entwicklungsumgebungen mit direktem Arbeitsbereichszugriff und Echtzeitvorschlägen.
  
  \item \textbf{Hybride Mensch-Agent-Workflows:} Optimale Arbeitsteilung zwischen Automatisierung und fachlicher Kontrolle sowie Werkzeuge für effiziente Delegation und Review.
  
  \item \textbf{Kosten-Nutzen-Optimierung:} Automatische Entscheidung wann Agent, wann Mensch basierend auf Aufgabenkomplexität, Deadline, Budgetbeschränkungen.
\end{itemize}

\subsection{Offene Forschungsfragen}

\begin{itemize}
  \item \textbf{Vertrauenswürdigkeit:} Wie lässt sich Vertrauenswürdigkeit messen oder absichern? Welche Rolle können formale Spezifikationen spielen?
  
  \item \textbf{Emergentes Verhalten:} Bei komplexen Multi-Agenten-Systemen: Wie kontrollieren/verstehen wir emergentes Verhalten?
  
  \item \textbf{Langzeitgedächtnis:} Wie skalieren semantische/episodische Gedächtnisse über Monate/Jahre? Forgetting vs. Retention Trade-offs?
  
  \item \textbf{Transfer Learning:} Können Agenten Wissen von Projekt A auf Projekt B transferieren? Wie lässt sich Domänenadaption für neue Codebasen gestalten?
\end{itemize}

\section{Schlusswort}

Die Arbeit zeigt, dass agentische Architekturen für Software-Engineering-Workflows praktisch umsetzbar sind und unter den gewählten Randbedingungen einen messbaren Nutzen entfalten können. Die entwickelte Referenzarchitektur, Implementierung und Evaluation bilden eine belastbare Grundlage für weiterführende Forschung und praktische Anwendungen.

Zentrale Erkenntnisse:

\begin{itemize}
  \item \textbf{Machbarkeit:} Agenten erreichen \pct{73} Erfolgsrate bei Refactoring \textendash{} ein belastbarer Hinweis auf praktische Einsetzbarkeit, jedoch keine hinreichende Evidenz für allgemeine Produktionsreife
  \item \textbf{Effizienz:} \pct{40} Token-Reduktion durch Optimierung \textendash{} Effizienzgewinne sind unter geeigneten Randbedingungen erreichbar
  \item \textbf{Sicherheit:} Defense-in-Depth erweist sich als zweckmäßiges Muster \textendash{} kontinuierliche Überprüfung bleibt jedoch erforderlich
  \item \textbf{Grenzen:} LLM-Halluzinationen, Kontextgrenzen, Latenz bleiben Herausforderungen
\end{itemize}

Die Technologie ist derzeit nicht hinreichend autonom für eine verantwortbare Vollautomatisierung, kann jedoch bereits als unterstützendes Werkzeug in Entwicklungsprozessen sinnvoll eingesetzt werden. Menschliche Freigabe bleibt sowohl technisch als auch ethisch unverzichtbar.

Zukünftige Arbeiten sollten nicht nur technische Verbesserungen fokussieren, sondern auch soziotechnische Fragen adressieren: Wie verändern Agenten Rollenbilder in der Softwareentwicklung? Wie lassen sich Übergänge in Organisationen verantwortlich gestalten? Und wie kann menschliche Expertise langfristig gesichert werden?

Die Kombination menschlicher Urteilsfähigkeit, Kreativität und Problemlösungskompetenz mit agentischer Automatisierung, Skalierung und Konsistenz kann zu einer verbesserten Unterstützung im Software Engineering beitragen \textendash{} sofern diese Systeme methodisch sauber evaluiert und verantwortungsvoll eingesetzt werden.

Konkrete Empfehlungen für Praktiker:

\begin{itemize}
  \item Beginnen Sie mit klar abgegrenzten, gut getesteten Teil-Workflows (z. B. Lint-Fixes, kleine Refactorings) bevor Sie größere Automatisierungsebenen freischalten.
  \item Implementieren Sie schrittweise menschliche Freigabepunkte für kritische Aktionen und messen Sie kontinuierlich Metriken wie Review-Akzeptanz und Regressionen.
  \item Nutzen Sie Vektor-basierte Retrieval-Mechanismen für langfristige Projekte, um Kontextkosten zu reduzieren, und automatisieren Sie Summarisierungspipelines für alte Commits.
  \item Planen Sie regelmäßige Sicherheits-Audits und erweitern Sie Test-Suites um adversariale Prompt-Tests.
\end{itemize}

Die vorgeschlagenen Schritte erleichtern den verantwortungsvollen Einsatz agentischer Systeme im Alltag und reduzieren Risiken beim Übergang in produktive Umgebungen.
